\section{Mechanism-Stratified Coupling Slopes}
\label{supp:slopes}

\textit{Purpose of this supplement.}\;
This supplement tabulates the five mechanism-stratified coupling
slopes (within-cell, schema-arm, encoder-arm, configuration-arm,
combined-cell) with cluster-robust standard errors and the
CI-width reportability gate applied in the main text.

The coupling-slope family complements the variance ratios with a
direct linear-regression view: each slope regresses
KB argmax accuracy on BioRED ex-NEG F1 at a different aggregation
level. All five slopes are reported descriptively because their
confidence intervals are too wide for a confirmatory directional
claim.

\subsection{Why a family rather than a pooled slope}

A pooled mixed-effects regression $\text{KB} \sim \text{BioRED} +
(1\,|\,\text{cell}) + (1\,|\,\text{seed})$ would identify its fixed
effect from within-cell variation, where seed-level BioRED has
standard deviation $\approx 0.03$ and seed-level KB has standard
deviation $\approx 0.16$. The seed-level slope is dominated by
independent noise sources (initialisation, dropout, sampler order)
and is uninformative about the cell-level relationship. A single
between-cell slope is no better, because schema-induced and
configuration-induced variation operate through different pathways
(label-space change versus optimiser dynamics on a fixed head) and
need not yield the same coupling. The family decomposition reports
each pathway separately.

\subsection{Slope estimates}

\begin{table}[H]
\centering
\caption{Coupling-slope estimates with 95\% cluster-bootstrap CIs.
The reportability gate (\S\ref{sec:s7-gate}) labels a slope
``inconclusive'' if its CI exceeds 0.30 in width.}
\label{tab:s7-slopes}
\small
\begin{tabular}{@{}lrrrl@{}}
\toprule
\textbf{Slope} & \textbf{$\hat\beta$} & \textbf{95\% CI} & \textbf{CI width} & \textbf{Label} \\
\midrule
$\beta_{\text{within-cell}}$       & $+1.10$  & $[+0.40, +1.75]$  & $1.35$  & inconclusive \\
$\beta_{\text{schema-arm}}$        & $+0.95$  & $[+0.69, +1.41]$  & $0.72$  & inconclusive \\
$\beta_{\text{encoder-arm}}$       & $+1.23$  & $[+0.73, +6.37]$  & $5.65$  & inconclusive \\
$\beta_{\text{configuration-arm}}$ & $-3.01$  & $[-13.00, +4.54]$ & $17.54$ & inconclusive \\
$\beta_{\text{combined-cell}}$     & $-0.91$  & not estimable     & ---     & --- \\
\bottomrule
\end{tabular}
\end{table}

The combined-cell slope's point estimate $-0.91$ is the unweighted
mean of the per-wave combined-cell slopes ($\hat\beta = +1.19$ in
the schema-comparison wave; $\hat\beta = -3.01$ in the
configuration wave); the wave-interaction CI is $[-9.13, +0.73]$
and does not exclude a large between-wave slope difference, so we
do not report a single pooled CI on
$\beta_{\text{combined-cell}}$.

\subsection{Reportability gate}
\label{sec:s7-gate}

The gate of 0.30 CI width matches the pre-committed slope-magnitude
partition's narrowest bin (``weak'', $|\beta| < 0.3$). A CI wider
than 0.30 cannot adjudicate the boundary between weak and moderate
coupling, so the slope receives no qualitative label.

\subsection{Why every slope is inconclusive}

The configuration-arm slope is fitted on 9 cell means, the schema-arm
and encoder-arm slopes on 12. These counts are below the threshold
at which OLS slopes have stable confidence intervals when between-cell
variation is large relative to within-cell variation. The
configuration arm illustrates the problem: 9 cells span a wide KB
range (0.150 to 0.777) but a narrow BioRED range (0.300 to 0.387),
and a few cells dominate the regression.

\subsection{What the family does support}

Despite the wide CIs, the point-estimate pattern is consistent with
the variance-ratio findings. The schema-arm and encoder-arm
estimates are positive and similar ($+0.95$ to $+1.23$), echoing the
schema-comparison wave's positive cross-metric correlations. The
configuration-arm estimate is negative and large in magnitude,
consistent with the variance-ratio reversal in that wave. The
within-cell estimate of $+1.10$ exceeds the noise-domination
expectation of near-zero, suggesting a small number of cells have
unusually high seed-level coupling rather than a stable population
property.

The variance ratio in \S\ref{sec:results-rq4} provides the more stable directional
claim: it is computed at the cell level over the same data but with
an operation less sensitive to between-cell range asymmetries.